% Created 2026-05-04 lun 16:57
% Intended LaTeX compiler: pdflatex
\documentclass[10pt,a4paper]{report}
\usepackage[utf8]{inputenc}
\usepackage{lmodern}
\usepackage[T1]{fontenc}
\usepackage[top=1in, bottom=1.in, left=1in, right=1in]{geometry}
\usepackage{graphicx}
\usepackage{longtable}
\usepackage{float}
\usepackage{wrapfig}
\usepackage{rotating}
\usepackage[normalem]{ulem}
\usepackage{amsmath}
\usepackage{textcomp}
\usepackage{marvosym}
\usepackage{wasysym}
\usepackage{amssymb}
\usepackage{amsmath}
\usepackage[theorems, skins]{tcolorbox}
\usepackage[version=3]{mhchem}
\usepackage[numbers,super,sort&compress]{natbib}
\usepackage{natmove}
\usepackage{url}
\usepackage[cache=false]{minted}
\usepackage[strings]{underscore}
\usepackage[linktocpage,pdfstartview=FitH,colorlinks,
linkcolor=blue,anchorcolor=blue,
citecolor=blue,filecolor=blue,menucolor=blue,urlcolor=blue]{hyperref}
\usepackage{attachfile}
\usepackage{setspace}
\author{Marcos Bujosa}
\date{\today}
\title{calcprop-export — exportación a AMC, Moodle y Jupyter}
\begin{document}

\tableofcontents

\part{Dando formato a las preguntas}
\label{sec:org6fb95f7}

\chapter{Estructura del paquete \texttt{calcprop-export}}
\label{sec:org877216f}

\begin{minted}[frame=lines,fontsize=\scriptsize,linenos=]{python}
from calcprop_quiz import *

<<Formato AMC>>
<<Formato AMC Verdadero/Falso>>
<<Formato AMC version Profe>>
<<Formato AMC con lastchoice>>
<<Formato AMC y multicolumna>>
<<Formato AMC y multicolumna version Profe>>
<<Formato AMC con lastchoice y multicolumna>>

<<Método codchar para codificar caracteres no ASCII en \LaTeX{}>>

<<Quiz para Moodle>>
<<Quiz para Moodle con valores>>
<<Quiz para Moodle versión Profe>>
<<Quiz para Moodle versión Profe con valores>>
<<Quiz VF para Moodle>>
<<Formato Moodle de preguntas de elección múltiple>>
<<Formato Moodle de preguntas de elección múltiple para el Profe>>

<<Quiz para Moodle con LastChoice>>
<<Quiz para Moodle con LastChoice con valores>>
<<Quiz VF para Moodle con LastChoice>>
<<Formato Moodle de preguntas de elección múltiple con LastChoice>>

<<Quiz para Jupyter>>

\end{minted}
\chapter{Preguntas en \LaTeX{} para usar con \href{https://www.auto-multiple-choice.net/}{Auto-Multiple-Choice}}
\label{sec:org9e98c1e}

Recorremos con un bucle la lista \texttt{cuestiones} de manera que en cada paso \texttt{c} es una nueva cuestión.
Si \texttt{c[2]} es cierta (si se verifica la precondición necesaria para formular la cuestión),
añadimos dicha cuestión en la cadena de cuestiones \texttt{s} del siguiente modo:
si la cuestión es verdadera (si \texttt{c[1]==True}) escribimos \texttt{\textbackslash{}correctchoice\{}; y si no, escribimos \texttt{\textbackslash{}wrongchoice\{}.
Seguidamente escribimos el enunciado de la cuestión, \texttt{c[0]}, cerramos el comando de \LaTeX{} y saltamos de línea de texto, \texttt{\}\textbackslash{}n}.
\begin{minted}[frame=lines,fontsize=\scriptsize,linenos=]{python}
for c in cuestiones:
    if c[2]:
        s = s + '       ' \
               + ('\\correctchoice{' if c[1] else '\\wrongchoice  {') + c[0] + '}\n'
\end{minted}

En el caso de las preguntas en versión para el profe, si \texttt{c[2]} no es cierta, añadimos dicha cuestión en la cadena \texttt{ex} de \emph{cuestiones excluidas}, seguida de un punto y coma y la proposición \texttt{c[1]} correspondiente a dicha cuestión.
Así, en el formato AMC para el profesor aparecerá, en la parte de explicación de cada pregunta, el listado de preguntas excluidas dado el enunciado de esa pregunta.

\begin{minted}[frame=lines,fontsize=\scriptsize,linenos=]{python}
for c in cuestiones:
    if c[2]:
        s = s + '       ' \
              + ('\\correctchoice{' if c[1] else '\\wrongchoice  {') + c[0] + '}\n'
    else:
        ex = ex + ' cuestion: ' + c[0] + '; ' + c[1] + '\n'
\end{minted}

Cada variante tendrá una opción correcta si también añadimos la opción: \emph{``ninguna de las anteriores es correcta''}.
El comando \texttt{\textbackslash{}lastchoices\textbackslash{}n} asegura que ésta sea la última.
Indicamos que es incorrecta si alguna de las opciones en la lista \texttt{cuestiones} es \texttt{True}, o que ésta es la opción correcta en caso contrario.
\begin{minted}[frame=lines,fontsize=\scriptsize,linenos=]{python}
s = s + '       ' + '\\lastchoices\n'
s = s + '       ' \
      + ('\\wrongchoice  {'if any([c[1] for c in cuestiones])else'\\correctchoice{')\
      + OpcPorDefecto + '}\n'
\end{minted}
\section{Variante para preguntas Verdadero/Falso}
\label{sec:org84b2359}

Las preguntas de ejercicios del tipo Verdadero/Falso no tienen pre-condición, por tanto \texttt{c[2]} no existe.
\begin{minted}[frame=lines,fontsize=\scriptsize,linenos=]{python}
for c in cuestiones:
    s = s + (" " * 7) + ( r'\correctchoice{' if c[1] else r'\wrongchoice{' ) + c[0] + '}\n'
\end{minted}
\begin{enumerate}
\item Formatos para \href{https://www.auto-multiple-choice.net/}{AMC}
\label{sec:org3719d77}

\begin{minted}[frame=lines,fontsize=\scriptsize,linenos=]{python}
def AMC (nombre, etiqueta, enunciado, cuestiones, opc=["","Ninguna de las anteriores"]):
    InstruccionesAux = opc[0]
    OpcPorDefecto    = opc[1]
    s = '\\element{' + nombre + '}{' + InstruccionesAux + '\n'
    s = s + ' \\begin{questionmult}{' + nombre + '-' + str(etiqueta) + '}\n'
    s = s + '  ' + enunciado + '\n'

    s = s + '     \\begin{choices}\n'
    for c in cuestiones:
        if c[2]:
            s = s + '       ' \
                   + ('\\correctchoice{' if c[1] else '\\wrongchoice  {') + c[0] + '}\n'
    s = s + '     \\end{choices}\n'

    s = s + ' \\end{questionmult} '
    s = s + '}\n\n'
    return s
\end{minted}

\begin{minted}[frame=lines,fontsize=\scriptsize,linenos=]{python}
def AMC_VF (nombre, etiqueta, enunciado, cuestiones, opc=["","Ninguna de las anteriores"]):
    InstruccionesAux = opc[0]
    OpcPorDefecto    = opc[1]
    s = '\\element{' + nombre + '}{' + InstruccionesAux + '\n'
    s = s + ' \\begin{questionmult}{' + nombre + '-' + str(etiqueta) + '}\n'
    s = s + '  ' + enunciado + '\n'

    s = s + '     \\begin{choices}\n'
    for c in cuestiones:
        s = s + (" " * 7) + ( r'\correctchoice{' if c[1] else r'\wrongchoice{' ) + c[0] + '}\n'
    s = s + '     \\end{choices}\n'

    s = s + ' \\end{questionmult} '
    s = s + '}\n\n'
    return s
\end{minted}



\begin{minted}[frame=lines,fontsize=\scriptsize,linenos=]{python}
def AMCProfe (nombre, etiqueta, enunciado, cuestiones, opc=["","Ninguna de las anteriores"]):
    InstruccionesAux = opc[0]
    OpcPorDefecto    = opc[1]
    ex = ''
    s = '\\element{' + nombre + '}{' + InstruccionesAux + '\n'
    s = s + ' \\begin{questionmult}{' + nombre + '-' + str(etiqueta) + '}\n'
    s = s + '  ' + enunciado + '\n'

    s = s + '     \\begin{choices}\n'
    for c in cuestiones:
        if c[2]:
            s = s + '       ' \
                  + ('\\correctchoice{' if c[1] else '\\wrongchoice  {') + c[0] + '}\n'
        else:
            ex = ex + ' cuestion: ' + c[0] + '; ' + c[1] + '\n'
    s = s + '     \\end{choices}\n'
    if ex:
        s = s + '   \\explain{' + ex + '            }\n'
    s = s + ' \\end{questionmult} '
    s = s + '}\n\n'
    return s
\end{minted}

\begin{minted}[frame=lines,fontsize=\scriptsize,linenos=]{python}
def AMClastCh (nombre, etiqueta, enunciado, cuestiones, opc=["","Ninguna de las anteriores"]):
    InstruccionesAux = opc[0]
    OpcPorDefecto    = opc[1]
    s = '\\element{' + nombre + '}{' + InstruccionesAux + '\n'
    s = s + ' \\begin{questionmult}{' + nombre + '-' + str(etiqueta) + '}\n'
    s = s + '  ' + enunciado + '\n'

    s = s + '     \\begin{choices}\n'
    for c in cuestiones:
        if c[2]:
            s = s + '       ' \
                   + ('\\correctchoice{' if c[1] else '\\wrongchoice  {') + c[0] + '}\n'
    s = s + '       ' + '\\lastchoices\n'
    s = s + '       ' \
          + ('\\wrongchoice  {'if any([c[1] for c in cuestiones])else'\\correctchoice{')\
          + OpcPorDefecto + '}\n'
    s = s + '     \\end{choices}\n'

    s = s + ' \\end{questionmult} '
    s = s + '}\n\n'
    return s
\end{minted}

\begin{minted}[frame=lines,fontsize=\scriptsize,linenos=]{python}
def AMCmc (nombre, etiqueta, enunciado, cuestiones, c, opc=["","Ninguna de las anteriores"]):
    InstruccionesAux = opc[0]
    OpcPorDefecto    = opc[1]
    s = '\\element{' + nombre + '}{' + InstruccionesAux + '\n'
    s = s + ' \\begin{questionmult}{' + nombre + '-' + str(etiqueta) + '}\n'
    s = s + '  ' + enunciado + '\n'
    s = s + '   \\begin{multicols}{' + str(c) + '}\\AMCBoxedAnswers\n'
    s = s + '     \\begin{choices}\n'
    for c in cuestiones:
        if c[2]:
            s = s + '       ' \
                   + ('\\correctchoice{' if c[1] else '\\wrongchoice  {') + c[0] + '}\n'
    s = s + '     \\end{choices}\n'
    s = s + '   \\end{multicols}\n'
    s = s + ' \\end{questionmult} '
    s = s + '}\n\n'
    return s
\end{minted}

\begin{minted}[frame=lines,fontsize=\scriptsize,linenos=]{python}
def AMCmcProfe (nombre, etiqueta, enunciado, cuestiones, c, opc=["","Ninguna de las anteriores"]):
    InstruccionesAux = opc[0]
    OpcPorDefecto    = opc[1]
    ex = ''
    s = '\\element{' + nombre + '}{' + InstruccionesAux + '\n'
    s = s + ' \\begin{questionmult}{' + nombre + '-' + str(etiqueta) + '}\n'
    s = s + '  ' + enunciado + '\n'
    s = s + '   \\begin{multicols}{' + str(c) + '}\\AMCBoxedAnswers\n'
    s = s + '     \\begin{choices}\n'
    for c in cuestiones:
        if c[2]:
            s = s + '       ' \
                  + ('\\correctchoice{' if c[1] else '\\wrongchoice  {') + c[0] + '}\n'
        else:
            ex = ex + ' cuestion: ' + c[0] + '; ' + c[1] + '\n'
    s = s + '     \\end{choices}\n'
    s = s + '   \\end{multicols}\n'
    if ex:
        s = s + '   \\explain{' + ex + '            }\n'
    s = s + ' \\end{questionmult} '
    s = s + '}\n\n'
    return s
\end{minted}

\begin{minted}[frame=lines,fontsize=\scriptsize,linenos=]{python}
def AMClastChmc (nombre, etiqueta, enunciado, cuestiones, c, opc=["","Ninguna de las anteriores"]):
    InstruccionesAux = opc[0]
    OpcPorDefecto    = opc[1]
    s = '\\element{' + nombre + '}{' + InstruccionesAux + '\n'
    s = s + ' \\begin{questionmult}{' + nombre + '-' + str(etiqueta) + '}\n'
    s = s + '  ' + enunciado + '\n'
    s = s + '   \\begin{multicols}{' + str(c) + '}\\AMCBoxedAnswers\n'
    s = s + '     \\begin{choices}\n'
    for c in cuestiones:
        if c[2]:
            s = s + '       ' \
                   + ('\\correctchoice{' if c[1] else '\\wrongchoice  {') + c[0] + '}\n'
    s = s + '       ' + '\\lastchoices\n'
    s = s + '       ' \
          + ('\\wrongchoice  {'if any([c[1] for c in cuestiones])else'\\correctchoice{')\
          + OpcPorDefecto + '}\n'
    s = s + '     \\end{choices}\n'
    s = s + '   \\end{multicols}\n'
    s = s + ' \\end{questionmult} '
    s = s + '}\n\n'
    return s
\end{minted}
\end{enumerate}
\section{Ejemplo de uso}
\label{sec:org1f135ce}

\begin{minted}[frame=lines,fontsize=\scriptsize,linenos=]{python}
from calcprop_export import *

enunciado = "Indique qué afirmaciones son verdaderas:"
banco = [
     ("Todo alumno que va a clase aprueba",   False),
     ("Todo alumno que suspende va a clase",  False),
     ("Todo alumno que sabe aprueba",         True ),
     ("Si aprueban todos eres buen profesor", False),
     ("Si suspenden todos eres mal profesor", False),
     ("Si suspenden todos es frustrante",     True ),]
GenVar = iter( ProblemaVF (enunciado, banco, 3) ) # tres preguntas por variante

nombre     = "EjemploVF"
directorio = "../ejemplos/"
with open(directorio + nombre + ".tex","w") as f:
    for i in range(4):                            # cuatro variantes
        var = (next(GenVar))
        f.write(AMC_VF(nombre,var[0],var[1],var[2]))
\end{minted}
\chapter{Preguntas en \texttt{xml} para usar con Moodle}
\label{sec:org89c0709}

Para generar los ficheros \texttt{xml} vamos a seguir un método
indirecto. El motivo es que escribir en \texttt{xml} puede ser tedioso
debido a los caracteres no ASCII y las expresiones matemáticas de
\LaTeX{}. La solución más práctica que he encontrado es generar desde
Python un fichero \LaTeX{} que a su vez emplea el paquete
\href{https://ctan.org/pkg/moodle}{moodle}. Así, al
compilar el fichero de \LaTeX{} con \texttt{pdflatex} obtendremos, tanto
una versión en pdf de las preguntas (¡Algo mucho más cómodo a la hora
de trabajar), como un fichero \texttt{xml} para ser importado en
Moodle. Además, este método indirecto nos evita generar la estructura
de los ficheros de preguntas en \texttt{xml} al estilo Moodle\dots{} !la
estructura en \LaTeX{} es mucho más sencilla! (o al menos, más
familiar).

\begin{minted}[frame=lines,fontsize=\scriptsize,linenos=]{python}
s = "\\documentclass[11pt]{article}\n\n"
s = s + "\\usepackage[cm,headings]{fullpage}\n\n"
s = s + "\\usepackage{moodle}\n\n"
s = s + "\\usepackage{graphicx}\n\n"
s = s + "\\usepackage{fancyvrb}\n\n"
s = s + "\\ifPDFTeX                     % FOR LATEX and PDFLATEX\n"
s = s + "    \\usepackage[utf8]{inputenc}   % necessary\n"
s = s + "    \\usepackage[OT1] {fontenc}    % necessary\n"
s = s + "\\else                         % assuming XELATEX or LUALATEX\n"
s = s + "    \\usepackage{fontspec}\n"
s = s + "\\fi\n\n"
s = s + auxLaTeX + "\n" + "\\newcommand\\peque{}" + "\n\n"
s = s + "\\begin{document}\n\n"
\end{minted}

Con las cuestiones de opción múltiple, el 100$\backslash$% de la nota de la
pregunta se obtiene si se marcan todas las correctas (indicadas con
\texttt{\textbackslash{}item*} y el 100$\backslash$% de la nota se pierde si se marcan todas la
incorrectas (las que aparecen con =\item =).

\textbf{Advertencia:} Lamentablemente Moodle no permite que la
puntuación global de ninguna pregunta de opción múltiple sea negativa
(es imposible que las preguntas de opción múltiple tengan esperanza
nula; en contraste con
\href{https://www.auto-multiple-choice.net/}{AMC}, donde se puede
programar cualquier regla de puntuación).

\begin{minted}[frame=lines,fontsize=\scriptsize,linenos=]{python}
for c in cuestiones:
    s = s + "       " + ( itemBuena(c[3]) if c[1] else itemMala(c[3]) ) + codchar(c[0]) + "\n"
\end{minted}

\begin{minted}[frame=lines,fontsize=\scriptsize,linenos=]{python}
for c in cuestiones:
    if c[2]:
        s = s + '       ' + (itemBuena(b,c[3]) if c[1] else itemMala(m,c[3])) + codchar(c[0]) + "\n"
    else:
        ex = ex + ' cuestion: ' + c[0] + '; ' + c[1] + '\n'
\end{minted}

Si ninguna respuesta es correcta, al compilar el fichero de \LaTeX{}
obtendremos un error (Moodle exige que al menos haya una respuesta
correcta). Por ello, podemos incluir la opción \emph\{``Las demás
  respuestas son incorrectas''\}. Lamentablemente, Moodle tampoco tiene
un comando del tipo \texttt{\textbackslash{}lastchoice}, que fuerza a que esta opción
aparezca en último lugar aunque se barajen las opciones; por eso pongo
\emph{las demás} en lugar de \emph{las anteriores} como en los exámenes
con \href{https://www.auto-multiple-choice.net/}{AMC}. Esta opción
será falsa si alguna de las demás es \texttt{True}.

\label{cuestion alternativa por defecto para Moodle}
\begin{minted}[frame=lines,fontsize=\scriptsize,linenos=]{python}
v = [c[1] for c in cuestiones].count(True)
cuestiones = cuestiones + [(codchar(lastchoice), (False if v else True), 1, '')]
\end{minted}
\section{\texttt{QuizMoodle}, \texttt{QuizMoodleLastCh}, \texttt{QuizVFMoodle} y \texttt{QuizVFMoodleLastCh}}
\label{sec:org63747e6}

\texttt{QuizMoodle} genera un archivo \LaTeX{} con las variantes de un \texttt{ProblemaTipo}.
Sus argumentos son el \texttt{nombre} del \texttt{ProblemaTipo}, el \texttt{directorio} donde queremos crear el fichero de \LaTeX{}, y la lista \texttt{problema} que contiene las variantes.
La cadena \texttt{s} contiene el preámbulo del archivo.
A la cadena \texttt{s} se le concatena el inicio del entorno \texttt{quiz}, cuyo primer argumento es el \texttt{nombre} del problema (del que escribiremos las distintas versiones).

\begin{itemize}
\item Con \texttt{with open()} creamos el archivo \texttt{nombre.tex} en el \texttt{directorio} que hayamos indicado; y entonces
\begin{itemize}
\item Escribimos la cadena \texttt{s} en dicho archivo.

\item Luego recorremos todas las \texttt{var=iantes de la lista =problema} con un \texttt{for}.
Por cada variante llamamos a \texttt{MoodleMulti} para escribir una pregunta de opción múltiple con los argumentos \texttt{nombre}, número de variante (\texttt{var[0]}), enunciado de la variante (\texttt{var[1]}), y las cuestiones de la variante (\texttt{var[2]}).

\item Una vez terminado el bucle \texttt{for}, escribimos el cierre del entorno \texttt{quiz} y el final del documento de \LaTeX{}.
\end{itemize}
\end{itemize}

\begin{minted}[frame=lines,fontsize=\scriptsize,linenos=]{python}
def creaDiccionario(x, key='key'):
    return x if isinstance(x, dict) else {key: x}
\end{minted}

\begin{minted}[frame=lines,fontsize=\scriptsize,linenos=]{python}
def QuizMoodle (nombre, directorio, problema, opc=["",""]):
    <<Método auxiliar que transforma un objeto en diccionario si no era ya un diccionario>>
    problema = creaDiccionario(problema, nombre)
    auxLaTeX = opc[0]
    <<Preámbulo fichero \LaTeX{} con paquete moodle>>
    s = s + "\\begin{quiz}{" + nombre + "}\n\n"
    with open(directorio + nombre + ".tex","w") as f:
        f.write(s)
        for i,nom in enumerate(problema):
            for var in problema[nom]:
                f.write( MoodleMulti(codchar(nom), var[0], codchar(var[1]), var[2]) )
        f.write("\\end{quiz}\n\n\\end{document}\n")

\end{minted}

\begin{minted}[frame=lines,fontsize=\scriptsize,linenos=]{python}
def QuizMoodleConVariables (nombre, directorio, problema, environment, Valores, opc=["",""]):
    <<Método auxiliar que transforma un objeto en diccionario si no era ya un diccionario>>
    problema = creaDiccionario(problema, nombre)
    auxLaTeX = opc[0]
    <<Preámbulo fichero \LaTeX{} con paquete moodle>>
    s = s + "\\begin{quiz}{" + nombre + "}\n\n"
    with open(directorio + nombre + ".tex","w") as f:
        f.write(s)
        for i,nom in enumerate(problema):
            for var in problema[nom]:
                template = environment.from_string(MoodleMulti(codchar(nom), var[0], codchar(var[1]), var[2]))
                content = template.render(Valores())
                f.write(content)
        f.write("\\end{quiz}\n\n\\end{document}\n")

\end{minted}

\begin{minted}[frame=lines,fontsize=\scriptsize,linenos=]{python}
def QuizMoodleProfe (nombre, directorio, problema, opc=["",""]):
    <<Método auxiliar que transforma un objeto en diccionario si no era ya un diccionario>>
    problema = creaDiccionario(problema, nombre)
    auxLaTeX = opc[0]
    <<Preámbulo fichero \LaTeX{} con paquete moodle>>
    s = s + "\\begin{quiz}{" + nombre + "}\n\n"
    with open(directorio + nombre + ".tex","w") as f:
        f.write(s)
        for i,nom in enumerate(problema):
            for var in problema[nom]:
                f.write( MoodleMultiProfe(codchar(nom), var[0], codchar(var[1]), var[2]) )
        f.write("\\end{quiz}\n\n\\end{document}\n")

\end{minted}

\begin{minted}[frame=lines,fontsize=\scriptsize,linenos=]{python}
def QuizMoodleProfeConVariables (nombre, directorio, problema, environment, Valores, opc=["",""]):
    <<Método auxiliar que transforma un objeto en diccionario si no era ya un diccionario>>
    problema = creaDiccionario(problema, nombre)
    auxLaTeX = opc[0]
    <<Preámbulo fichero \LaTeX{} con paquete moodle>>
    s = s + "\\begin{quiz}{" + nombre + "}\n\n"
    with open(directorio + nombre + ".tex","w") as f:
        f.write(s)
        for i,nom in enumerate(problema):
            for var in problema[nom]:
                template = environment.from_string(MoodleMultiProfe(codchar(nom), var[0], codchar(var[1]), var[2]))
                content = template.render(Valores())
                f.write(content)
        f.write("\\end{quiz}\n\n\\end{document}\n")

\end{minted}

Si necesitamos que las variantes añadan la \hyperref[cuestion alternativa por defecto para Moodle]{cuestion alternativa por defecto para Moodle}
para así asegurarnos de que siempre hay una respuesta correcta,
entonces empleamos el método \texttt{QuizMoodleLastCh}, que hace lo
mismo que el anterior pero llamando a \texttt{MoodleMultiLastCh}.

\begin{minted}[frame=lines,fontsize=\scriptsize,linenos=]{python}
def QuizMoodleLastCh (nombre, directorio, problema, opc=["","Las demás opciones son falsas"]):
    <<Método auxiliar que transforma un objeto en diccionario si no era ya un diccionario>>
    auxLaTeX      = opc[0]
    OpcPorDefecto = opc[1]
    <<Preámbulo fichero \LaTeX{} con paquete moodle>>
    s = s + "\\begin{quiz}{" + nombre + "}\n\n"
    with open(directorio + nombre + ".tex","w") as f:
        f.write(s)
        for i,nom in enumerate(creaDiccionario(problema, nombre)):
            for var in problema[nom]:
                f.write(MoodleMultiLastCh(codchar(nom),var[0],codchar(var[1]),var[2],opc[1]))
        f.write("\\end{quiz}\n\n\\end{document}\n")

\end{minted}

\begin{minted}[frame=lines,fontsize=\scriptsize,linenos=]{python}
def QuizMoodleLastChConVariables (nombre, directorio, problema, environment, Valores, opc=["","Las demás opciones son falsas"]):
    <<Método auxiliar que transforma un objeto en diccionario si no era ya un diccionario>>
    auxLaTeX      = opc[0]
    OpcPorDefecto = opc[1]
    <<Preámbulo fichero \LaTeX{} con paquete moodle>>
    s = s + "\\begin{quiz}{" + nombre + "}\n\n"
    with open(directorio + nombre + ".tex","w") as f:
        f.write(s)
        for i,nom in enumerate(creaDiccionario(problema, nombre)):
            for var in problema[nom]:
                template = environment.from_string(MoodleMultiLastCh(codchar(nom), var[0], codchar(var[1]), var[2]))
                content = template.render(Valores())
                f.write(content)
        f.write("\\end{quiz}\n\n\\end{document}\n")

\end{minted}

\texttt{QuizVFMoodle} genera un archivo \LaTeX{} con las variantes de un \texttt{ProblemaVF}.
Los argumentos son el \texttt{nombre} del \texttt{ProblemaVF}, el \texttt{directorio} donde crear el archivo, el iterador \texttt{GenVar} que genera las variantes, y el \texttt{num}-ero de variantes que se desea.
\begin{minted}[frame=lines,fontsize=\scriptsize,linenos=]{python}
def QuizVFMoodle (nombre, directorio, GenVar, num, opc=["",""]):
    auxLaTeX      = opc[0]
    <<Preámbulo fichero \LaTeX{} con paquete moodle>>
    s = s + "\\begin{quiz}{" + nombre + "}\n\n"
    with open(directorio + nombre + ".tex","w") as f:
        f.write(s)
        for i in range(num):
            var = next(GenVar)
            f.write( MoodleMulti (codchar(nombre), var[0], codchar(var[1]), var[2]) )
        f.write("\\end{quiz}\n\n\\end{document}\n")
\end{minted}

Cuando necesitamos que las variantes añadan la \hyperref[cuestion alternativa por defecto para Moodle]{cuestion alternativa por defecto para Moodle}, empleamos el método \texttt{QuizVFMoodleLastCh}, que hace lo mismo que el anterior pero llamando a \texttt{MoodleMultiLastCh}.
\begin{minted}[frame=lines,fontsize=\scriptsize,linenos=]{python}
def QuizVFMoodleLastCh (nombre, directorio, GenVar, num, opc=["","Las demás opciones son falsas"]):
    auxLaTeX      = opc[0]
    <<Preámbulo fichero \LaTeX{} con paquete moodle>>
    s = s + "\\begin{quiz}{" + nombre + "}\n\n"
    with open(directorio + nombre + ".tex","w") as f:
        f.write(s)
        for i in range(num):
            var = next(GenVar)
            f.write(MoodleMultiLastCh(codchar(nombre),var[0],codchar(var[1]),var[2],opc[1]))
        f.write("\\end{quiz}\n\n\\end{document}\n")
\end{minted}
\section{\texttt{MoodleMulti} y \texttt{MoodleMultiLastCh}}
\label{sec:orgcabcd75}

\texttt{MoodleMulti} genera una cadena de caracteres con los comandos \LaTeX{} correspondientes a una pregunta de opción múltiple.
Los parámetros del método \texttt{MoodleMulti} son: \texttt{nombre} (el nombre del problema) y \texttt{variante} (el número de la variante)\footnote{Que separados por un guión son empleados como nombre de cada variante del problema de opción múltiple dentro del código \LaTeX{}.}, el \texttt{enunciado} de la variante y la lista de \texttt{cuestiones} de la variante.
\begin{minted}[frame=lines,fontsize=\scriptsize,linenos=]{python}
def MoodleMulti (nombre, variante, enunciado, cuestiones):
    <<Escritura del entorno multi>>
\end{minted}

El método \texttt{MoodleMulti} escribe el código \LaTeX{} de un entorno \texttt{multi} del paquete \href{https://ctan.org/pkg/moodle}{moodle} con los siguientes parámetros opcionales para dicho entorno: \texttt{multiple} para que sea una pregunta de opción múltiple y \texttt{points} el número de opciones de respuesta de la pregunta (la longitud de \texttt{cuestiones}).

\begin{minted}[frame=lines,fontsize=\scriptsize,linenos=]{python}
<<Escritura de cada item de cuestión en Moodle>>

ex = ''
s = " \\begin{multi}[multiple, points=" + str(len(cuestiones)-1) +"]"
s = s + "{" + codchar(nombre) + "-" + str(variante) + "}\n"
s = s + "    " + enunciado + "\n"

<<Escritura de las cuestiones para Moodle>>

s = s + " \\end{multi}\n\n"
return s
\end{minted}

\begin{minted}[frame=lines,fontsize=\scriptsize,linenos=]{python}
def MoodleMultiProfe (nombre, variante, enunciado, cuestiones):
    <<Escritura aclaración en Moodle>>
    <<Escritura del entorno multi para el Profe>>
\end{minted}

\begin{minted}[frame=lines,fontsize=\scriptsize,linenos=]{python}
def feedback(texto):
    return r",feedback={"+codchar(texto)+"}"
\end{minted}

\begin{minted}[frame=lines,fontsize=\scriptsize,linenos=]{python}
<<Conteo de respuestas correctas e incorrectas>>
<<Escritura de cada item de cuestión en Moodle para Profe>>

ex = ''
s = " \\begin{multi}[multiple, fractiontol=5.1"  +"]"
s = s + "{" + codchar(nombre) + "-" + str(variante) + "}\n"
s = s + "    " + enunciado + "\n"

<<Escritura de las cuestiones para Moodle para el Profe>>

s = s + " \\end{multi}\n\n"
return s
\end{minted}

\begin{minted}[frame=lines,fontsize=\scriptsize,linenos=]{python}
b = 0; m = 0;
for c in cuestiones:
    print(c)
    if c[2]:
        if c[1]:
            b+=1
        else:
            m+=1
\end{minted}

\begin{minted}[frame=lines,fontsize=\scriptsize,linenos=]{python}
def itemBuena(aclaracion):
    return ("\\item[feedback={" + codchar(aclaracion) + "}]* ") if aclaracion else r"\item* "

def itemMala(aclaracion):
    return ("\\item[feedback={" + codchar(aclaracion) + "}]  ") if aclaracion else r"\item  "
\end{minted}

\begin{minted}[frame=lines,fontsize=\scriptsize,linenos=]{python}
def itemBuena(numBuenas,aclaracion):
    return ("\\item[fraction=" + str( round(100/numBuenas) ) + feedback(aclaracion) + "]") if numBuenas else "\\item*"

def itemMala(numMalas,aclaracion):
    return ("\\item[fraction=" + str(-round(100/numMalas) ) + feedback(aclaracion) + "]") if numMalas else "\\item"
\end{minted}

Cuando requerimos añadir la \hyperref[cuestion alternativa por defecto para Moodle]{cuestion alternativa por defecto para Moodle} llamamos a \texttt{MoodleMultiLastCh} que añade la opción \emph{``Las demás opciones son falsas''} a la lista \texttt{cuestiones}.
\begin{minted}[frame=lines,fontsize=\scriptsize,linenos=]{python}
def MoodleMultiLastCh (nombre, variante, enunciado, cuestiones, \
                       lastchoice="Las demás opciones son falsas"):
    <<cuestion alternativa por defecto para Moodle>>
    <<Escritura del entorno multi>>

\end{minted}
\subsection{Codificación de caracteres no \texttt{ASCII}}
\label{sec:org083ee02}

\begin{minted}[frame=lines,fontsize=\scriptsize,linenos=]{python}
def codchar(s):
     s = s.replace("á", "\\'{a}")
     s = s.replace("é", "\\'{e}")
     s = s.replace("í", "\\'{i}")
     s = s.replace("ó", "\\'{o}")
     s = s.replace("ú", "\\'{u}")
     s = s.replace("ñ", "\\~{n}")
     return s
\end{minted}
\section{Ejemplo de uso de \texttt{QuizMoodleLastCh} y \texttt{QuizVFMoodleLastCh}}
\label{sec:org10a0e76}

\begin{minted}[frame=lines,fontsize=\scriptsize,linenos=]{python}
from calcprop_export import *
p = ProblemaTipo( \
      [   "Considere ",
          [
              Supuesto("$\\mathcal{A}$ ", v("A")),
              Supuesto("$\\mathcal{B}$ ", v("B")),
          ],
          [
              Supuesto("y que $\\mathcal{A}\\Rightarrow\\mathcal{C}$. ",    \
                       v("A") >> v("C")),
              Supuesto("y que $\\mathcal{B}\\Leftrightarrow\\mathcal{D}$. ",\
                       v("B") ** v("D")),
          ],
          "Indique qué opción es correcta: ",
          [
              Cuestion("Entonces $\\mathcal{C}$ es verdadero", v("C")),
              Cuestion("Entonces $\\mathcal{D}$ es verdadero", v("D")),
              Cuestion("Entonces $\\mathcal{E}$ es verdadero", v("E"), v("D")),
          ],
      ])
nombre     = "EjemploMoodle"
directorio = "../ejemplos/"
QuizMoodleLastCh (nombre, directorio, p)

\end{minted}


\begin{minted}[frame=lines,fontsize=\scriptsize,linenos=]{python}
from calcprop_export import *

enunciado = "Indique qué afirmaciones son verdaderas:"
banco = [
  ("Todo alumno que va a clase aprueba",   False),
  ("Todo alumno que suspende va a clase",  False),
  ("Todo alumno que sabe aprueba",         True ),
  ("Si aprueban todos eres buen profesor", False),
  ("Si suspenden todos eres mal profesor", False),
  ("Si suspenden todos es frustrante",     True ),]

b = [(codchar(p[0]),p[1]) for p in banco]

GenVar = iter( ProblemaVF (codchar(enunciado), b, 3) )

nombre     = "EjemploVFMoodle"
directorio = "../ejemplos/"

QuizVFMoodleLastCh (nombre, directorio, GenVar, 4)

\end{minted}
\end{document}
